\documentclass[fleqn,10pt]{olplainarticle}
% Use option lineno for line numbers 

\title{Conditional Probabilistic Anatomical Reasoning for Tumor Segmentation with Multimodal Priors}

\author[1]{First Author}
\author[2]{Second Author}
\affil[1]{Address of first author}
\affil[2]{Address of second author}

\keywords{Keyword1, Keyword2, Keyword3}

\begin{abstract}
Please provide an abstract of no more than 300 words. Your abstract should explain the main contributions of your article, and should not contain any material that is not included in the main text. 
\end{abstract}

\begin{document}

\flushbottom
\maketitle
\thispagestyle{empty}

\section{Introduction}

\section{Methodology}
\subsection{Conditional Probabilistic Anatomical Reasoning}
\label{sec:conditional_reasoning}
\begin{figure}
    \centering
    \includegraphics[width=1\linewidth]{figs/PGM.png}
    \caption{ Anatomical Reasoning Graph.}
    \label{fig:pgm}
\end{figure}

We formulate tumor segmentation as a conditional
probabilistic anatomical reasoning problem that
jointly infers tumors and their host anatomical
structures by integrating image-dependent
segmentation priors with anatomical and
logical constraints.
Let $X$ denote the observed medical image,
$Z$ denote the tumor segmentation, and $H$
denote the host anatomical segmentation.
We further denote the cancer type by $C$
and the data-agnostic anatomical knowledge
base by $K$.
Our framework consists of three complementary
components: (1) observation consistency,
(2) tumor reasoning, and (3) host reasoning.

\paragraph{Image-Dependent Segmentation Priors.}

We construct four complementary segmentation
priors to provide image-dependent evidence
for tumor and host anatomical reasoning.

For tumor reasoning, we employ cancer-type
prompting to obtain 2D and 3D tumor
segmentation priors from BP and VT,
respectively:
\begin{equation}
\label{eq:tumor_priors}
\begin{aligned}
B_C &= f_{\mathrm{BP}}(X,C),\\
V_C &= f_{\mathrm{VT}}(X,C),
\end{aligned}
\end{equation}
where $B_C$ provides slice-wise tumor
segmentation evidence and $V_C$ provides
volumetric tumor segmentation evidence.

For host reasoning, we obtain two complementary
anatomical segmentation priors.
First, we leverage a data-agnostic anatomical
knowledge base $K_C$ to construct anatomical
prompts for VT, yielding a cancer-specific
anatomical segmentation prior:
\begin{equation}
\label{eq:kb_prior}
V_K = f_{\mathrm{VT}}(X,K_C).
\end{equation}
Second, we employ TotalSegmentator to
generate a cancer-agnostic anatomical segmentation
prior without explicit prompting:
\begin{equation}
\label{eq:ts_prior}
T = f_{\mathrm{TS}}(X).
\end{equation}
Although the anatomical knowledge base itself
is independent of the observed image, the
corresponding VT segmentation output $V_K$
is image-dependent.
Consequently, the framework incorporates
two tumor segmentation priors,
$\{B_C,V_C\}$, and two host anatomical
segmentation priors, $\{V_K,T\}$.

\paragraph{Conditional Probabilistic Formulation.}

Given the observed image, cancer type,
and anatomical knowledge base, we model
the joint conditional distribution of
tumor and host segmentations as
\begin{equation}
\label{eq:conditional_joint}
\boxed{
p_\theta(Z,H\mid X,C,K)
=
\frac{1}{\mathcal{Z}_\theta(X,C,K)}
\phi_{\mathrm{obs}}
\phi_{\mathrm{tumor}}
\phi_{\mathrm{host}}
}
\end{equation}
where $\mathcal{Z}_\theta(X,C,K)$ is the
normalizing partition function.
The three nonnegative potentials are defined as
\begin{equation}
\label{eq:reasoning_potentials}
\begin{aligned}
\phi_{\mathrm{obs}}
&=
\phi_{\mathrm{obs}}(Z,H;X),\\
\phi_{\mathrm{tumor}}
&=
\phi_{\mathrm{tumor}}(Z,H;B_C,V_C),\\
\phi_{\mathrm{host}}
&=
\phi_{\mathrm{host}}(H;V_K,T).
\end{aligned}
\end{equation}
This conditional formulation directly
incorporates the image-dependent segmentation
priors without assuming that they are
independent of the observed image or
constitute data-independent Bayesian priors.

\paragraph{Observation Consistency.}

The observation consistency potential
evaluates whether a hypothesized tumor--host
configuration is compatible with the
observed imaging evidence. We employ statistical rejection to suppress
candidate configurations that are inconsistent
with the observed image, \textbf{while develop statistical recall to discover more candidates that are consistent with the accepted tumor characteristics}.
Specifically, we define
\begin{equation}
\label{eq:observation_potential}
\phi_{\mathrm{obs}}(Z,H;X)
=
\exp\left[
-\lambda_{\mathrm{obs}}
E_{\mathrm{obs}}(Z,H;X)
\right],
\end{equation}
where $E_{\mathrm{obs}}$ is a
statistical inconsistency score (Section~\ref{sec:reasoning_details}) and
$\lambda_{\mathrm{obs}}$ controls the
strength of observation consistency.
Candidate configurations with greater
statistical inconsistency receive lower
compatibility scores, thereby reducing
their contribution to the joint conditional
distribution.

\paragraph{Tumor Reasoning.}

The tumor reasoning potential integrates
the cancer-type-prompted 2D and 3D tumor
segmentation priors with the inferred
host anatomical structures.
We formulate tumor reasoning through
set-theoretic operations and logical
constraints:
\begin{equation}
\label{eq:tumor_reasoning}
\begin{aligned}
\phi_{\mathrm{tumor}}
&(Z,H;B_C,V_C) =
\phi_{\mathrm{BP}}(Z;B_C)
\phi_{\mathrm{VT}}^{\mathrm{tumor}}(Z;V_C)
\phi_{\mathrm{logic}}(Z,H).
\end{aligned}
\end{equation}
The 2D segmentation potential
$\phi_{\mathrm{BP}}$ evaluates the
consistency between the candidate tumor
segmentation and slice-wise tumor evidence.
The 3D segmentation potential
$\phi_{\mathrm{VT}}^{\mathrm{tumor}}$
evaluates consistency with volumetric
tumor evidence.
The logical potential
$\phi_{\mathrm{logic}}$ encodes
anatomical relationships between tumors
and their host structures.
In particular, set-theoretic relations
such as containment, intersection,
and exclusion are used to evaluate
the anatomical plausibility of candidate
tumor configurations.
This enables complementary 2D and 3D
segmentation evidence to be integrated
under explicit anatomical constraints.

\paragraph{Host Reasoning.}

The host reasoning potential integrates
two complementary sources of image-dependent
anatomical evidence: the KB-prompted VT
segmentation prior and the non-prompted
TotalSegmentator segmentation prior.
We formulate the host reasoning potential as
\begin{equation}
\label{eq:host_reasoning}
\begin{aligned}
\phi_{\mathrm{host}}(H;V_K,T)
={}\phi_{\mathrm{VT}}^{\mathrm{anat}}(H;V_K) \cdot
\phi_{\mathrm{TS}}(H;T).
\end{aligned}
\end{equation}
The KB-prompted VT potential incorporates
patient-specific anatomical evidence
obtained through knowledge-guided
segmentation, whereas the TotalSegmentator
potential provides complementary anatomical
evidence without explicit prompting.
A cancer sits in its primary site with high prior but may have invaded neighbours or spread, so the host is a set with weights, not one organ.

\paragraph{Joint Tumor--Host Inference.}

The complete conditional reasoning model
can be expanded as
\begin{equation}
\label{eq:expanded_reasoning}
\boxed{
\begin{aligned}
p_\theta(Z,H\mid X,C,K)
\propto\;&\phi_{\mathrm{obs}}(Z,H;X) \\
&\cdot\phi_{\mathrm{BP}}(Z;B_C) \cdot\phi_{\mathrm{VT}}^{\mathrm{tumor}}
(Z;V_C) \cdot\phi_{\mathrm{logic}}(Z,H) \\
&\cdot\phi_{\mathrm{VT}}^{\mathrm{anat}}
(H;V_K) \cdot\phi_{\mathrm{TS}}(H;T).
\end{aligned}
}
\end{equation}
The final tumor and host segmentations
are estimated through maximum a posteriori
inference:
\begin{equation}
\label{eq:map_inference}
(\hat Z,\hat H)
=
\arg\max_{Z,H}
p_\theta(Z,H\mid X,C,K).
\end{equation}
Equivalently, the inference objective
can be expressed in the log domain:
\begin{equation}
\label{eq:log_inference}
\begin{aligned}
(\hat Z,\hat H)
=
\arg\max_{Z,H}\;
\log\phi_{\mathrm{obs}} +\log\phi_{\mathrm{tumor}} +\log\phi_{\mathrm{host}}.
\end{aligned}
\end{equation}
By integrating complementary tumor
segmentation priors, knowledge-guided
and non-prompted anatomical priors,
and explicit anatomical constraints,
the proposed framework enables
conditional probabilistic reasoning
over tumor--host configurations
under uncertainty.

\subsection{Pan-cancer knowledge base and prompt planning}
\label{sec:kb_planning}
% Drafted 4 Oct 2026 from knowledge base v0.5.6 (knowledge_base/*.yaml) and the planner (pancia_kb/planner.py,
% adapter.py). No equation of Section~\ref{sec:conditional_reasoning} is changed; this section defines how $K$ is
% built and how the prompt plan $K_C$ of Eq.~\eqref{eq:kb_prior} is produced for each scan.

This section describes how the anatomical knowledge base $K$ is built once for all cancer types, and how, for each scan,
it is turned into the final VT prompt plan $K_C$ of Eq.~\eqref{eq:kb_prior}. $K$ is a curated, image-independent graph;
$K_C$ is a finite list of prompt jobs, and $V_K=f_{\mathrm{VT}}(X,K_C)$ is the stack of masks VT returns for these jobs.
No component of $K$ or of the planner uses patient sex, and no rule or constant is specific to an organ, a cancer type
or a dataset.

\paragraph{Anatomical entities.}
$K$ contains 191 anatomical entities. Each entity has a category (solid organ, hollow organ, gland, bone, vessel,
lymph-node station, muscle, fascia, fat, duct, space, nerve, landmark or region) and a laterality (left, right, midline
or bilateral). It also records how it can be observed: the TotalSegmentator classes that represent it for CT and for
MR, if any; a VT main phrase with up to two aliases, each graded as in-vocabulary, near-vocabulary or
out-of-vocabulary against the published VT training vocabulary; its craniocaudal span between vertebral anchors;
whether it is always present, variably present or absent after common surgery; a typical volume range; and, for CT, a
typical Hounsfield-unit range. Finally, each entity that can host a tumor carries a tumor phrase (for example
\emph{liver tumor}). Thirty-four entities are \emph{anchors}: structures that TotalSegmentator segments reliably and
that therefore fix the coordinate frame of a scan (vertebrae, major vessels and large organs). Landmarks are tied to
vertebral levels, so that every spatial prior in $K$ is expressed relative to the patient's own skeleton rather than in
millimetres.

\paragraph{Anatomical relations.}
Entities are connected by 401 typed, directed relations. Ten relation kinds are defined, of which nine are used:
\emph{adjacent\_to} (156), \emph{has\_part} (94), \emph{drains\_lymph\_to} (57), \emph{lies\_in} (23),
\emph{supplied\_by} (17), \emph{invested\_by} (16), \emph{drained\_by} (16), \emph{landmark\_for} (13) and
\emph{has\_duct} (9); the tenth, \emph{wall\_of}, is reserved. Each relation kind has a fixed spatial gate and a default
priority. The gate states where the target may lie relative to the source: inside it (\emph{has\_part}), in a shell
around it (\emph{invested\_by}), in a neighbouring box or directional shell (\emph{adjacent\_to}), in a corridor along
a vessel or duct (\emph{supplied\_by}, \emph{drained\_by}, \emph{has\_duct}) or in a body region
(\emph{drains\_lymph\_to}); \emph{lies\_in} and \emph{landmark\_for} carry no gate. Gates are later used to check VT
outputs against structures that are already segmented, and priorities order prompts when the budget is exceeded. The relations encode the anatomy that staging systems refer to: direct extension
into a neighbouring organ, invasion of an investing fascia or fat, and lymphatic drainage to a regional station.

\paragraph{Cancer-type spread sets.}
For each cancer type $C$, $K$ stores a weighted host set $\mathcal{H}_C$, which is the host set used in
Section~\ref{sec:reasoning_details}. Only two items are curated per cancer type: the primary organ (prior $1.0$) and the
usual distant metastatic sites (prior $0.15$). Organs at risk of local invasion (prior $0.3$) are not curated; they are
derived automatically as the entities reached from the primary organ by an \emph{adjacent\_to} or \emph{invested\_by}
relation. Each cancer type also has a short ensemble of tumor phrases (for example \emph{lung cancer}, \emph{lung
tumor}, \emph{lung mass}), and generic templates produce phrases for any host without a curated one. The current base
covers 18 cancer types. Adding a new cancer type therefore requires only its primary organ, its metastatic sites and,
optionally, its phrase ensemble; all other knowledge is shared.

\paragraph{Planning $K_C$ for a scan.}
For a given scan, the planner builds $K_C$ deterministically in a single pass, without calling any segmentation model.
Its inputs are the modality, the cancer type $C$, a summary of the TotalSegmentator output $T$ (which classes are
present, their volumes and their extents), and the initial tumor priors $B_C$ and $V_C$ of Eq.~\eqref{eq:tumor_priors}.
It proceeds in five steps.

\emph{(i) Frame.} The vertebrae found in $T$ define a craniocaudal ruler. The landmarks of $K$ are placed on this ruler,
and every entity span is converted to a slice range of the scan. This step also determines which entities lie inside the
field of view.

\emph{(ii) Anchors.} Each anchor is labelled \emph{found} (present in $T$ with a plausible volume and position),
\emph{expected} (inside the field of view but missing from $T$) or \emph{missing} (outside the field of view). Found
anchors give the reference structures against which the gates of later prompts are evaluated.

\emph{(iii) Tumor evidence.} A component of $B_C$ or $V_C$ is \emph{validated} when the two priors agree on it. For each
host $h\in\mathcal{H}_C$ the planner records an evidence level $e_h\in\{0,1,2\}$: $2$ if a validated component lies in
or within a reach of $10$\,mm of $h$, $1$ if a component from one prior does, and $0$ otherwise. It also records the
side of the primary tumor for paired organs. The planning weight of the host is $w_h(1+e_h)$, where $w_h$ is its
spread-set prior. The initial tumor priors therefore influence only the order and the extent of the anatomical
prompts; they never remove a host from $\mathcal{H}_C$.

\emph{(iv) Prompt tiers.} Prompts are generated in four tiers. Tier~0 adds ruler prompts (spine, sacrum, hip bone) only
when $T$ provides neither vertebrae nor fallback landmarks, so that a frame can always be placed. Tier~1 completes the anchors that are expected but missing from $T$, each gated by its
spatial prior. Tier~2 adds \emph{anatomical profiles} by walking the relation graph: two hops from the primary organ,
one hop from each secondary host, and the structures that can be reached from any host with tumor evidence
$e_h>0$. Tier~T contains the tumor prompts: one prompt with the tumor phrase of each host of $\mathcal{H}_C$ inside the
field of view. Tier~T is issued for every such host regardless of the evidence $e_h$, so that a lesion missed by both
initial priors can still be proposed. For the same reason, when the initial VT tumor mask is empty or smaller than
$0.5$\,ml, the cancer-type phrase ensemble is added for the primary organ. This is a wording ensemble for the same
model, not a second rater: the initial and prompted VT outputs are counted as one model in tumor reasoning.

\emph{(v) Budget.} The plan is limited to $80$ prompts. Tiers~0, 1 and T are protected; only tier~2 is trimmed, lowest relation
priority first, with items derived from a host given a one-step priority boost.

The output $K_C$ is the list of prompt jobs that survive the budget. Each job records the entity, its side, the phrases to
issue (the main phrase and its first two aliases), its tier, its gate, and, for tumor prompts, the host and the host
prior $w_h$. VT is run once per scan on the full volume with all jobs of $K_C$; there is no re-planning after inference.
For each job, the masks of its three phrases are fused by voxel-wise majority (at least two of three). The gate is then
applied as an output filter: each connected component of the fused mask is kept only when at least $80\%$ of its
voxels lie inside the region allowed by the gate. A completion structure that is empty after filtering is recorded as
absent when $K$ marks it as variable or removable by surgery, and as not found otherwise. The kept organ masks form $V_K$, which enters host reasoning
through $\phi^{\mathrm{anat}}_{\mathrm{VT}}$. The masks of the tier-T jobs are tumor proposals; they are added to the VT
tumor prior $V_C$ as candidate parts, each tagged with its prompted host, and are subject to the same host, logic and
observation reasoning as every other part.

Two properties of this construction matter for the reasoning that follows. First, although $K$ is independent of the
image, $K_C$ is not: it depends on $T$ through the frame and the anchors, and on the initial tumor priors through the
evidence levels. This dependence is restricted to which structures are prompted and in what order, so it changes the
coverage of $V_K$ but not the meaning of any mask. Second, because tier~T prompts every plausible host whether or not it
contains a tumor, VT is regularly asked to segment a tumor that is absent. Its output for such a prompt can be a false
positive inside the prompted host or elsewhere in the scan. These proposals are therefore treated as priors only, and
whether they are accepted is decided by the image evidence $\phi_{\mathrm{obs}}$ of
Section~\ref{sec:reasoning_details}.

\subsection{Reasoning details}
\label{sec:reasoning_details}
% Drafted 4 Oct 2026 from the implemented ledger (pancia_kb/ledger.py) and the agreed design of statistical recall and
% rejection (to be implemented as ledger v4). All constants are fixed a priori and identical for every cancer type.

This section specifies each potential of Eq.~\eqref{eq:expanded_reasoning}. All operations are performed per scan on the
native voxel grid of $X$, on which $B_C$, $V_C$, $V_K$ and $T$ are defined. Distances $d(\cdot,\cdot)$ are Euclidean in
millimetres. The \emph{acquisition plane} is orthogonal to the axis with the largest voxel spacing when the ratio of
largest to smallest spacing is at least $1.5$, and is the axial plane otherwise; all slice-wise operations use this plane.
Every constant is fixed before inference, shared by all cancer types, and not tuned on clinical endpoints.

\textit{Variables.} Tumor reasoning (below) produces a finite set of candidate parts $\mathcal{A}=\{a_1,\dots,a_n\}$, which
are disjoint voxel sets. Each part has a binary label $z_a\in\{0,1\}$ and a host assignment $h(a)$, and the tumor
segmentation is $Z=\bigcup_{a:\,z_a=1}a$. The host segmentation $H=\{H_h\}_{h\in\mathcal{H}_C}$ assigns one region to each
host of the cancer type. Given $H$, the potentials of Eq.~\eqref{eq:expanded_reasoning} factorise over the parts, so that
\begin{equation}
\label{eq:part_posterior}
\log p_\theta(Z\mid H,X,C,K) = \sum_{a\in\mathcal{A}} z_a\Bigl[\underbrace{\log\phi_{\mathrm{logic}}(a,H)}_{0\ \text{or}\ -\infty}
+\underbrace{\operatorname{logit}\pi_a}_{\phi_{\mathrm{BP}}\,\phi^{\mathrm{tumor}}_{\mathrm{VT}}}
+\underbrace{\log\phi_{\mathrm{obs}}(a,H;X)}_{\text{image evidence}}\Bigr] + \mathrm{const},
\end{equation}
and the MAP solution of Eq.~\eqref{eq:map_inference} sets $z_a=1$ exactly when the bracket is non-negative. The three
paragraphs below define $\phi_{\mathrm{host}}$, the three tumor terms and $\phi_{\mathrm{obs}}$.

%--------------------------------------------------------------------------------------------------------------------
\paragraph{Host reasoning: $\phi_{\mathrm{TS}}(H;T)$ and $\phi^{\mathrm{anat}}_{\mathrm{VT}}(H;V_K)$.}
The knowledge base defines for cancer type $C$ a weighted host set $\mathcal{H}_C=\{(h,\kappa_h,w_h)\}$: the primary organ
($\kappa_h=$ primary, $w_h=1.0$), organs reached by local invasion ($\kappa_h=$ local, $w_h=0.3$) and metastatic sites
($\kappa_h=$ distant, $w_h=0.15$), with laterality where applicable. For each host, $T_h$ denotes the TotalSegmentator
labels mapped to $h$ by $K$, and $V_{K,h}$ the VT segmentation of the knowledge-base organ prompt for $h$ (voxel-wise
majority over the main phrase and its aliases). The two host potentials are containment constraints with a parsimony
term,
\begin{equation}
\label{eq:host_potentials}
\phi_{\mathrm{TS}}(H;T)=\prod_{h}\mathbf{1}\bigl[T_h\subseteq H_h\bigr]\,e^{-\epsilon|H_h|},\qquad
\phi^{\mathrm{anat}}_{\mathrm{VT}}(H;V_K)=\prod_{h}\mathbf{1}\bigl[V_{K,h}\subseteq H_h\bigr],
\end{equation}
with $\epsilon>0$ arbitrarily small. Their product is maximised by
\begin{equation}
\label{eq:host_map}
\hat H_h = T_h\cup V_{K,h},
\end{equation}
so each organ is supported by either prior, and $r_h=\mathrm{Dice}(T_h,V_{K,h})$ records the agreement of the two priors.

Special cases:
\begin{itemize}
\item When neither prior delineates the primary organ, a coarse knowledge-base region anchored on vertebral levels replaces
$\hat H_h$ and is flagged; a secondary host without any organ mask has no region.
\item When the planner finds the primary organ outside the field of view, no part can be assigned to it.
\end{itemize}

Derived regions:
\begin{itemize}
\item The host envelope is $E_h=\{v: d(v,\hat H_h)\le r_h\}$, with $r_h=10$\,mm for the primary and $3$\,mm for secondary hosts.
\item The contact zone is $E^{c}_h=\{v: d(v,\hat H_h)\le3\,\mathrm{mm}\}$.
\item The other-organ region $\bar O$ is the union of all TotalSegmentator structures except the primary organ, its parts and
its satellites, each filled slice-wise for enclosed holes. TotalSegmentator excludes tumor tissue from organ masks, so
without filling, a lesion enclosed by an organ would not count as inside it.
\end{itemize}

The host weights $w_h$ rank admitted lesions (the index lesion is a primary lesion) and do not enter acceptance.

%--------------------------------------------------------------------------------------------------------------------
\paragraph{Tumor reasoning: $\phi_{\mathrm{BP}}(Z;B_C)$, $\phi^{\mathrm{tumor}}_{\mathrm{VT}}(Z;V_C)$ and
$\phi_{\mathrm{logic}}(Z,H)$.}
Tumor reasoning (i)~builds the candidate parts $\mathcal{A}$ from the two tumor priors, (ii)~assigns each part a prior
probability from the support of $B_C$ and $V_C$, and (iii)~removes anatomically or logically infeasible parts.

\textit{(i) Candidate parts.}
\begin{itemize}
\item \emph{Claims.} Connected components (26-connectivity) of $B_C$, of $V_C$ and of every VT tumor prompt for a host
$h\in\mathcal{H}_C$ are claims if they are at least $0.5$\,ml. Inside the primary envelope, VT fragments closer than
$5$\,mm are grouped first and the floor is $0.2$\,ml.
\item \emph{Lesions.} Let $\mathcal{B}$ and $\mathcal{V}$ be the unions of the BP and VT claims. Lesions are anchored on the
volumetric prior. On every acquisition slice that contains $\mathcal{V}$, a 2D component of $\mathcal{B}\setminus\bar O$
that touches $\mathcal{V}$ (within one pixel) joins it with its full in-plane extent, and BP voxels within $5$\,mm of
$\mathcal{V}$ in 3D also join. The union is grouped at $5$\,mm. The remaining BP voxels form BP-only lesions.
\item \emph{Parts.} Each lesion $L$ is split into parts:
  \begin{itemize}
  \item the agreed part $L\cap\mathcal{B}\cap\mathcal{V}$;
  \item VT-only parts $L\cap\mathcal{V}\setminus\mathcal{B}$, separated by whether BP proposes anything on that acquisition
  slice within the host envelope;
  \item BP-only parts $L\cap\mathcal{B}\setminus\mathcal{V}$, split once by the sign of the voxel log-likelihood ratio
  (Eq.~\eqref{eq:llr}) and separated by whether VT proposes anything within $10$\,mm.
  \end{itemize}
  Connected parts below $0.5$\,ml are not scored and follow their lesion.
\item \emph{Recall parts.} Unproposed voxels ($\notin\mathcal{B}\cup\mathcal{V}$) within $10$\,mm of an accepted lesion,
inside its host envelope and outside $\bar O$, that have a positive voxel log-likelihood ratio and form connected parts of
at least $0.5$\,ml touching the lesion.
\end{itemize}

\textit{(ii) Tumor-prior potentials.} For each part $a$, let $\beta_a$ and $\nu_a$ be the states of BP and VT on $a$:
\emph{proposes} (the part lies in the model's claims), \emph{silent} (the model proposes nothing there: for BP, no claim on
the part's acquisition slices within the host envelope; for VT, no claim within $10$\,mm), or \emph{active} (the model
proposes tissue there but not this part). The two tumor potentials act jointly as a Bernoulli prior on the part labels,
\begin{equation}
\label{eq:tumor_prior}
\phi_{\mathrm{BP}}(Z;B_C)\,\phi^{\mathrm{tumor}}_{\mathrm{VT}}(Z;V_C)=\prod_{a\in\mathcal{A}}\pi_a^{z_a}(1-\pi_a)^{1-z_a},
\qquad
\pi_a=\begin{cases}
0.9 & \beta_a=\nu_a=\text{proposes},\\
0.75 & \text{one proposes, the other silent},\\
0.5 & \text{one proposes, the other active},\\
0.25 & \text{neither proposes (recall part)}.
\end{cases}
\end{equation}
Silence of one model is missing evidence, not evidence against. The value $0.75$ lies halfway, in log-odds, between
agreement and contradiction.

\textit{(iii) Logical potential.} $\phi_{\mathrm{logic}}$ is a product of hard constraints,
\begin{equation}
\label{eq:logic}
\phi_{\mathrm{logic}}(Z,H)=\prod_{a\in\mathcal{A}}\Bigl(\prod_{k}\Gamma_k(a,H)\Bigr)^{z_a},\qquad\Gamma_k\in\{0,1\},
\end{equation}
so a part violating any constraint has $\log\phi_{\mathrm{logic}}=-\infty$ and cannot be accepted. The constraints are:
\begin{itemize}
\item[$\Gamma_1$] \emph{Prompt gate.} A part from a VT prompt for host $h$ lies at least $50\%$ inside $E_h$.
\item[$\Gamma_2$] \emph{Shape.} The claim is neither a slab nor organ-shaped. A slab has median in-plane coverage of the
host $\ge0.4$ and a through-plane to in-plane extent ratio $\le0.25$; it is exempt when the other model claims the same
voxels. An organ-shaped claim has $\ge60\%$ of the host volume, median slice-wise Dice $\ge0.6$ and $\ge80\%$ inside
$\hat H_h$; it is excluded when the host priors agree ($r_h\ge0.7$), and otherwise kept on the primary host with weight
$0.3$.
\item[$\Gamma_3$] \emph{Assignment and containment.} The part's lesion is assigned to $h(a)\in\mathcal{H}_C$ and
satisfies $|L\cap\bar O|<0.5\,|L|$. Assignment is to the primary if $\ge50\%$ of $L$ lies in the primary envelope or $L$
touches $E^{c}_{\text{primary}}$; otherwise to the secondary host whose envelope holds $\ge50\%$ of $L$.
\item[$\Gamma_4$] \emph{Feasibility.} A primary lesion touches $E^{c}_h$, and the primary organ is in the field of view.
\item[$\Gamma_5$] \emph{Multifocality.} A second or further single-support primary lesion is a nodule: $\ge80\%$ inside
$\hat H_h$ (or attached), longest in-plane diameter $\ge10$\,mm, and extent $\ge3$ slices and $\ge10$\,mm.
\item[$\Gamma_6$] \emph{Invasion and spread.} A local-invasion lesion is contiguous with an accepted primary lesion
(within $3$\,mm). A non-contiguous lesion in an organ that is also a metastatic site is treated as distant, and otherwise
excluded. A distant lesion's prompt is coherent: $\ge80\%$ of the prompted mask lies inside $E_h$.
\end{itemize}
Rules that depend only on whether the second model agrees are not constraints: they are expressed by $\pi_a$ in
Eq.~\eqref{eq:tumor_prior} and decided by the image. These rules are "unconfirmed", "single model without VT in the
host", and "distant lesion without BP".

%--------------------------------------------------------------------------------------------------------------------
\paragraph{Statistical recall and rejection: $\phi_{\mathrm{obs}}(Z,H;X)$.}
The observation potential grounds every part in the observed image, including parts on which both tumor priors agree. It
is written as
\begin{equation}
\label{eq:obs_def}
\phi_{\mathrm{obs}}(Z,H;X)=\prod_{a\in\mathcal{A}}\Bigl(\mathbf{1}\bigl[\mathrm{BF}^{(1)}_a\ge\tau\bigr]\;
\mathbf{1}\bigl[\mathrm{BF}^{(2)}_a\ge\tau\bigr]\;\mathrm{BF}^{(2)}_a\Bigr)^{z_a},\qquad \tau=3,\ \lambda_{\mathrm{obs}}=1,
\end{equation}
i.e. $E_{\mathrm{obs}}=-\sum_a z_a\log\mathrm{BF}^{(2)}_a$ on parts that pass both evidence thresholds, and
$E_{\mathrm{obs}}=+\infty$ otherwise. $\tau=3$ is the conventional threshold for positive evidence~\cite{kass1995bayes}.
$\mathrm{BF}^{(1)}_a$ is the one-sided Bayes factor (is $a$ different from its local anatomy?), and $\mathrm{BF}^{(2)}_a$
the two-sided Bayes factor (does $a$ look like the accepted tumor rather than the local anatomy?).

\textit{Image features.}
\begin{itemize}
\item For voxel $v$, $\mathbf{f}(v)=(x_v,\mu_v,\sigma_v)$: the intensity, and the mean and standard deviation in its
$3\times3$ neighbourhood in the acquisition plane.
\item Within each scan, every feature is mapped to its empirical cumulative distribution over the pooled reference voxels.
Comparisons are made within the scan and are invariant to monotone intensity transformations, so the same procedure
applies to CT and MR across scanners and protocols.
\end{itemize}

\textit{References.}
\begin{itemize}
\item The local anatomy $A(a)$ is the set of voxels within $20$\,mm of $a$ that belong to $\hat H_{h(a)}$ or to any
TotalSegmentator structure, after removing all candidate parts dilated by $2$\,mm and eroding by one voxel. It contains the
host parenchyma and every surrounding structure: vessels, mediastinum, neighbouring solid organs, and the wall and content
of hollow organs.
\item The tumor reference $R_T$ is formed from the agreed parts that passed the one-sided gate, in the same host (the whole
scan if the host has fewer than $200$ such voxels), eroded by one voxel. If the scan has none, the single-support parts
that passed the gate are used instead. $R_T$ is therefore always validated by the image, and $a$ itself is always left out
($R_T\setminus a$).
\item $\hat p_T$ and $\hat p_A$ are Gaussian kernel density estimates (Scott bandwidth, at most $5000$ voxels) of
$\mathbf{f}$ on $R_T\setminus a$ and on $A(a)$.
\end{itemize}

\textit{Statistics.}
\begin{equation}
\label{eq:llr}
D(a)=\mathcal{E}\bigl(\mathbf{f}(a),\,\mathbf{f}(A(a))\bigr),\qquad
s(a)=\frac{1}{|a|}\sum_{v\in a}\ell(v),\quad
\ell(v)=\log\frac{\hat p_T(\mathbf{f}(v))}{\hat p_A(\mathbf{f}(v))},
\end{equation}
where $\mathcal{E}$ is the energy distance between the two feature distributions and $\ell(v)$ is the voxel log-likelihood
ratio.

\textit{Calibration.}
\begin{itemize}
\item Both statistics are calibrated within the scan against $M=200$ pseudo-parts: copies of $a$'s mask, or spatially
contiguous blocks of the same volume, placed at random inside a reference region.
\item Pseudo-parts reproduce the size, shape and spatial correlation of $a$, so no independence between voxels is assumed.
\item \emph{One-sided.} Pseudo-parts placed in $A(a)$ give the null distribution of $D$ and
$p_a=\bigl(1+\#\{m: D_m\ge D(a)\}\bigr)/(M+1)$. This is converted into the Bayes-factor bound of~\cite{sellke2001calibration},
\begin{equation}
\label{eq:bf1}
\mathrm{BF}^{(1)}_a=\frac{1}{-e\,p_a\ln p_a}\quad(p_a<e^{-1}),\qquad \mathrm{BF}^{(1)}_a=1\ \text{otherwise}.
\end{equation}
\item \emph{Two-sided.} Pseudo-parts drawn from $R_T\setminus a$ and from $A(a)$ give the sampling densities $g_T$ and $g_A$
of $s$, and
\begin{equation}
\label{eq:bf2}
\mathrm{BF}^{(2)}_a=\frac{g_T\bigl(s(a)\bigr)}{g_A\bigl(s(a)\bigr)}.
\end{equation}
\end{itemize}

\textit{Decision.} Substituting Eqs.~\eqref{eq:tumor_prior}, \eqref{eq:logic} and~\eqref{eq:obs_def} into
Eq.~\eqref{eq:part_posterior}, the MAP label of a feasible part ($\prod_k\Gamma_k=1$) is
\begin{equation}
\label{eq:decision}
\hat z_a=1\iff \mathrm{BF}^{(1)}_a\ge3,\quad \mathrm{BF}^{(2)}_a\ge3,\quad
\operatorname{logit}\pi_a+\log\mathrm{BF}^{(2)}_a\ge0.
\end{equation}
Because $\pi_a\ge0.25$ (prior odds $\ge1/3$), the last condition holds whenever $\mathrm{BF}^{(2)}_a\ge3$. Positive image
evidence is therefore the single acceptance threshold for every part. A part without image evidence is rejected even if
both models propose it, and a part with evidence is accepted even if one model proposes it or neither does. The priors set
the order of inference and decide which parts may form $R_T$; they are ablated in Section~X.

\textit{Inference order.} $R_T$ depends on which parts are accepted, so the MAP solution is computed in a fixed order.
\begin{enumerate}
\item \emph{Gate.} Every feasible lesion part is tested one-sided ($\mathrm{BF}^{(1)}\ge3$). Parts without evidence are
rejected.
\item \emph{Tumor reference.} $R_T$ is built from the gated parts as defined above.
\item \emph{Two-sided decision.} Every gated part is accepted or rejected by Eq.~\eqref{eq:decision} (statistical
rejection). A lesion that the logical admission would have dropped only for lack of a second model is admitted if its core
is accepted (statistical recovery). Single-support parts of accepted lesions that fail are removed (trimming).
\item \emph{Recall.} Recall parts of accepted lesions are tested by Eq.~\eqref{eq:decision} with $\pi=0.25$, and accepted
parts are added to their lesion (statistical recall).
\end{enumerate}
Accepted parts are re-assembled into lesions and re-ranked by $\phi_{\mathrm{logic}}$ and $w_h$ (primary first). The index
lesion is then measured in the acquisition plane.


\section*{Acknowledgments}

Additional information can be given in the template, such as to not include funder information in the acknowledgments section.

\bibliography{sample}

\end{document}